%% kapa: exam=1 round=20 year=2009 period=1 form=A subject=civil_law subject_ko=민법 seq=18 number=18
%% source: Q-Net 감정평가사 기출문제 (cst00309) artlSeq=5024178 file=extracted/1008/5024178/2009년 제20회 감정평가사 1차시험 기출문제/1교시.hwp
%% conversion: kordoc 4.24.51 -> tools/gen1.py (auto); review=proofread-sonnet; flags=none
\begin{question}{18}
점유자와 회복자의 관계에 있어서 옳은 내용만을 모두 고른 것은? (다툼이 있으면 판례에 의함)
\begin{framedtext}
ㄱ. 점유자는 선의로 점유한 것으로 추정되지만, 권원 없는 점유였음이 밝혀지면 곧 그 동안의 점유에 대한 선의의 추정은 깨어진다. \par
ㄴ. 선의의 점유자란 과실수취권을 포함한 본권을 가지고 있다고 오신하고 그 오신에 대하여 오신할 만한 정당한 근거가 있는 경우를 말한다. \par
ㄷ. 과실을 수취한 점유자는 그가 지출한 통상의 필요비의 상환을 회복자에게 청구할 수 있다. \par
ㄹ. 점유자가 점유물을 개량하기 위하여 지출한 금액 기타 유익비에 관하여는 그 가액의 증가가 현존한 경우에 한하여 회복자의 선택에 좇아 그 지출금액이나 증가액의 상환을 청구할 수 있다. \par
\end{framedtext}
\begin{choices}
  \item ㄱ, ㄴ
  \item ㄴ, ㄷ
  \item ㄴ, ㄹ
  \item ㄱ, ㄹ
  \item ㄷ, ㄹ
\end{choices}
\end{question}
\answer{3}
